\documentclass[border=10pt]{standalone}
\usepackage[no-math]{fontspec}
\usepackage[brazil]{babel}
\usepackage{amsmath,amssymb}
\IfFileExists{C:/Windows/Fonts/Lato_Regular.ttf}{%
  \setmainfont{Lato}[
    Path = C:/Windows/Fonts/,
    UprightFont = Lato_Regular.ttf,
    BoldFont = Lato_Bold.ttf,
    ItalicFont = Lato_Italic.ttf,
    BoldItalicFont = Lato_BoldItalic.ttf
  ]
  \setsansfont{Lato}[
    Path = C:/Windows/Fonts/,
    UprightFont = Lato_Regular.ttf,
    BoldFont = Lato_Bold.ttf,
    ItalicFont = Lato_Italic.ttf,
    BoldItalicFont = Lato_BoldItalic.ttf
  ]
}{%
  \RequirePackage[default]{lato}
}
\renewcommand{\familydefault}{\sfdefault}

\usepackage{xcolor}
\usepackage{tikz}
\usepackage{pgfplots}
\pgfplotsset{compat=1.18}
\usetikzlibrary{arrows.meta,positioning,shapes.geometric,fit,backgrounds,calc,decorations.pathreplacing,patterns}

% Paleta institucional UEFS / Bioengenharia de Solos
\definecolor{bsNavy}{HTML}{1F4268}
\definecolor{bsNavyDark}{HTML}{152E4A}
\definecolor{bsGold}{HTML}{B38A3D}
\definecolor{bsGoldDark}{HTML}{8C6927}
\definecolor{bsTeal}{HTML}{0E7C7B}
\definecolor{bsForest}{HTML}{2D6A4F}
\definecolor{bsAlert}{HTML}{9E2A2B}
\definecolor{bsEarth}{HTML}{8C5A3C}
\definecolor{bsWater}{HTML}{2E86AB}
\definecolor{bsInk}{HTML}{26313D}
\definecolor{bsMuted}{HTML}{5B6573}

\definecolor{bgNavy}{HTML}{EDF2F7}
\definecolor{bgGold}{HTML}{F8F3E9}
\definecolor{bgTeal}{HTML}{EBF6F5}
\definecolor{bgForest}{HTML}{EEF5F0}
\definecolor{bgAlert}{HTML}{FDF2F2}

\tikzset{
  >=Stealth,
  bsBox/.style={
    draw=#1, fill=white, rounded corners=5pt, line width=1.15pt,
    inner sep=7pt, align=center, font=\footnotesize\sffamily, text=bsInk
  },
  bsTint/.style={
    draw=#1, fill=#1!8!white, rounded corners=5pt, line width=1.15pt,
    inner sep=7pt, align=center, font=\footnotesize\sffamily, text=bsInk
  },
  bsBadge/.style={
    fill=#1, text=white, rounded corners=9pt, inner xsep=8pt, inner ysep=3pt,
    font=\scriptsize\bfseries\sffamily
  },
  bsTitle/.style={font=\normalsize\bfseries\sffamily, text=bsNavy, align=center},
  bsSub/.style={font=\footnotesize\sffamily, text=bsMuted, align=center},
  bsArrow/.style={->, >=Stealth, line width=1.35pt},
  bsDArrow/.style={<->, >=Stealth, line width=1.35pt}
}


\begin{document}
\begin{tikzpicture}
  \hyphenpenalty=10000\exhyphenpenalty=10000

  % Título
  \node[bsTitle] at (0, 3.25)
    {Integração Modelagem Hidráulica (HEC-RAS) e Restauração de APPs Fluviais};
  \node[bsSub] at (0, 2.80)
    {Fluxo de trabalho baseado em tensão cisalhante admissível para dimensionamento de margens};

  % Estágio 1
  \node[bsTint=bsNavy, text width=3.45cm, minimum height=3.25cm, anchor=north] (s1) at (-6.0, 2.15) {
    \vspace{0.32cm}
    \textbf{\color{bsNavy}Topobatimetria\\\& Hidrologia}\\[5pt]
    \scriptsize
    \textbullet\ MDT de alta resolução e seções do canal\\[2pt]
    \textbullet\ Hidrogramas de projeto:\\
    \hspace{0.2cm}$Q_{10},\; Q_{50},\; Q_{100}$\\[2pt]
    \textbullet\ Rugosidade inicial ($n$)
  };
  \node[bsBadge=bsNavy] at (s1.north) {ETAPA 1};

  % Estágio 2
  \node[bsTint=bsWater, text width=3.45cm, minimum height=3.25cm, anchor=north] (s2) at (-2.0, 2.15) {
    \vspace{0.32cm}
    \textbf{\color{bsWater}Simulação 1D/2D\\(HEC-RAS)}\\[5pt]
    \scriptsize
    \textbullet\ Perfis de linha d'água e planície de inundação\\[2pt]
    \textbullet\ Mapas de velocidade ($V$) por trecho\\[2pt]
    \textbullet\ Tensão cisalhante ($\tau_0$)
  };
  \node[bsBadge=bsWater] at (s2.north) {ETAPA 2};

  % Estágio 3
  \node[bsTint=bsAlert, text width=3.45cm, minimum height=3.25cm, anchor=north] (s3) at (2.0, 2.15) {
    \vspace{0.32cm}
    \textbf{\color{bsAlert}Seleção de\\Bioengenharia}\\[5pt]
    \scriptsize
    \textbullet\ Critério de estabilidade:\\
    \hspace{0.2cm}$\tau_0 > \tau_{\text{adm}}$\\[2pt]
    \textbullet\ Definição de enrocamento, gabião vivo ou feixes\\[2pt]
    \textbullet\ Fator de segurança marginal
  };
  \node[bsBadge=bsAlert] at (s3.north) {ETAPA 3};

  % Estágio 4
  \node[bsTint=bsForest, text width=3.45cm, minimum height=3.25cm, anchor=north] (s4) at (6.0, 2.15) {
    \vspace{0.32cm}
    \textbf{\color{bsForest}APP \&\\Monitoramento}\\[5pt]
    \scriptsize
    \textbullet\ Recomposição ciliar por zona de inundação\\[2pt]
    \textbullet\ Inspeção pós-cheia e batimetria de controle\\[2pt]
    \textbullet\ Calibração de rugosidade
  };
  \node[bsBadge=bsForest] at (s4.north) {ETAPA 4};

  % Setas diretas
  \draw[bsArrow, bsNavy]  (s1.east) -- (s2.west);
  \draw[bsArrow, bsWater] (s2.east) -- (s3.west);
  \draw[bsArrow, bsAlert] (s3.east) -- (s4.west);

  % Loop de retroalimentação inferior (com folga generosa abaixo das caixas)
  \draw[bsArrow, bsTeal, dashed, rounded corners=6pt]
    (s4.south) -- (6.0, -1.80) -- (-2.0, -1.80) -- (s2.south);
  \node[fill=bgTeal, draw=bsTeal, rounded corners=4pt, inner xsep=8pt, inner ysep=4pt, font=\scriptsize\bfseries, text=bsTeal]
    at (2.0, -1.80) {Retroalimentação pós-cheia: atualização de $n$ de Manning com o crescimento da mata ciliar};

\end{tikzpicture}
\end{document}
